\section{Where the Optimisations Do and Do Not Apply}
\label{sec:generality}

nx-oxcaml shaped the generalised analysis of Section~\ref{sec:counted}, so it
is a development workload, not an independent test. Two workloads the passes
were not adjusted for follow: a standard kernel suite, which gains, and a
cryptography code base, which does not. A port of the passes to stock OCaml
comes first.

\subsection{The Passes on Stock OCaml 5.3.0}
\label{sec:stock}
None of the passes depends on flambda2, so I ported them to OCaml 5.3.0.
Stock OCaml polls every loop, so with default flags only strength reduction
fires: the passes gain $+12.2\%$ (21/21) on microgpt. Counted
loops add $+4.9\%$ (17/21), for
$+17.6\%$ (21/21) in all over stock OCaml, and the
inner loop goes from 30 to 17 instructions per 4 multiply--adds. The best OxCaml build stays $+12.8\%$ (20/21)
ahead of the best OCaml build. OCaml's test suite (1,462 tests) gives the same
result test by test with and without the patches, also with counted loops
forced on.

\subsection{A Held-Out Kernel Suite: PolyBench}
\label{sec:polybench}
PolyBench/C~\cite{polybench} is a standard suite of loop-nest kernels
(linear algebra, stencils, data mining) written by others for compiler
research. I ported 17 of its kernels to plain OCaml as an ordinary programmer
would write them (flat row-major \code{float array}s, \code{for} loops,
\code{-unsafe}) and to C, and ran them on the pull-request compiler in a
separate 21-round run. All builds produce bit-identical results, and all
84 loops in the timed kernels are counted (Table~\ref{tab:polybench}). The
geometric-mean gain is $+19.7\%$ with array sizes fixed at compile time
(gemm $+108\%$) and $+7.5\%$ with sizes known only at run time. The
rules of Table~\ref{tab:applicability} predict that kernels whose inner loops
index arrays as \code{base + m*counter + c} should gain, and kernels limited by
memory bandwidth or a chain of dependent arithmetic should not.

\begin{table}[h]
  \centering
  \small
  \begin{tabular}{@{}lrrr@{}}
    \toprule
    kernel & counted loops vs passes off & inner-loop instructions & gap to scalar C \\
    \midrule
    gemm & $+108.2\%$ (21/21) & 19 $\to$ 8 & $2.24\times$ \\
    jacobi-2d & $+49.6\%$ (21/21) & 21 $\to$ 11 & $1.43\times$ \\
    fdtd-2d & $+48.7\%$ (21/21) & 18--22 $\to$ 10--12 & $1.34\times$ \\
    syrk & $+42.6\%$ (21/21) & 20 $\to$ 11 & $2.34\times$ \\
    doitgen & $+33.5\%$ (21/21) & 21 $\to$ 11 & $3.16\times$ \\
    floyd-warshall & $+26.2\%$ (21/21) & 23 $\to$ 15 & $2.12\times$ \\
    heat-3d & $+16.6\%$ (21/21) & 38 $\to$ 31 & $1.49\times$ \\
    atax & $+15.0\%$ (20/21) & 15 $\to$ 7 & $2.26\times$ \\
    2mm & $+10.5\%$ (20/21) & 19 $\to$ 11 & $3.65\times$ \\
    8 others & $-0.7$ to $+6.4\%$ & & $1.18$--$2.84\times$ \\
    \midrule
    geometric mean, 17 kernels & $+19.7\%$ & & $2.03\times$ \\
    \quad same 17 kernels, sizes known only at run time & $+7.5\%$ & & \\
    \bottomrule
  \end{tabular}
  \caption{PolyBench in plain OCaml with \texttt{-unsafe}: the passes with
  counted loops against the passes off, polls on in both (paired median,
  rounds won; sizes fixed at compile time unless noted). Two
  byte-identical builds differed by $-0.5$ to $+2.1\%$ per kernel. Others:
  gemver, trmm, mvt, gesummv, bicg, seidel-2d, covariance, correlation.}
  \label{tab:polybench}
\end{table}

The prediction holds. Gains are large where an inner loop indexes several
two-dimensional arrays this way (gemm, the stencils, syrk, doitgen) and the
passes remove most of its instructions, which counted loops allow by moving
the poll to the exit. Gains are small where something else
limits the loop: matrix--vector kernels on 32--64\,MB matrices that wait on
memory, column-strided access, and seidel-2d, whose loop shrinks from 37 to 15
instructions for only $+2\%$ because each update waits for the previous one.

Two further results bound the claim. With sizes known only at run time the
gain is $+7.5\%$: constant sizes let the compiler fold each row stride into
the loop,
while run-time strides stay as multiplications inside it and add register
pressure, so the passes remove a smaller share of its work (gemm's
loop shrinks from 30 to 27 instructions instead of 19 to 8). And OCaml stays
$2.03\times$ behind scalar C, because GCC keeps array-element accumulators in
registers and unrolls and interleaves the gemm loop, which no OCaml pass does.

\subsection{The oxcaml-crypto-study Repository}
\label{sec:crypto}
oxcaml-crypto-study~\cite{crypto_study} asks whether OxCaml can implement
cryptographic primitives natively, at Rust-like speed, instead of binding to C.
Its one primitive, ChaCha20~\cite{rfc8439}, exists as a boxed plain-OCaml
version and an OxCaml version with unboxed \texttt{int32\#} arrays,
stack-allocated state and \code{[@zero_alloc]}. I compiled both with the
pull-request compiler in eight configurations (the polling modes, each with
and without the pointer passes) and timed encryption of 1\,MiB\,+\,37 bytes
over 15 interleaved rounds; all passed the RFC~8439 test vector.

\subsection{Result: No Measurable Effect}
In every polling mode the assembly is \emph{byte-identical} with the four
passes on or off; the flags only drop back-edge polls in nine loops. The best
configuration measured $\times1.019$ (12/15 rounds), while two byte-identical
builds differed by 1.7\%, for four reasons:

\begin{enumerate}
  \item \textbf{Byte offsets, not array elements.} The hot XOR loop's
        \code{Bytes.get_int32_le plaintext (j*64 + i*4)} untags a byte offset
        with a \emph{right} shift and adds it unscaled; the passes match
        addresses scaled by a left shift.
  \item \textbf{The invariant is not a variable.} flambda2 inlines
        \code{j * 64} into the address as \code{(<< j 6)}; writing
        \code{let base = j * 64} gave byte-identical assembly.
  \item \textbf{The other loops are already optimal or not GC-free.} The copy
        and add loops index by the bare counter, the round loop calls a
        non-inlined function, and the boxed version allocates on every store.
  \item \textbf{The ceiling is low.} Each block makes 80 calls to a
        94-instruction quarter-round, about 7,500 instructions, against
        roughly 600 in loops the passes could touch.
\end{enumerate}

\subsection{An Applicability Characterisation}
Table~\ref{tab:applicability} summarises the results: the passes recover
C-style pointer walking over arrays, helping code whose cost is index
arithmetic and doing nothing for code whose cost is the computation itself.

\begin{table}[h]
  \centering
  \small
  \begin{tabularx}{\textwidth}{@{}lXX@{}}
    \toprule
    property & benefits & does not benefit \\
    \midrule
    addressing & array index \code{base + m*counter + c}, scaled by a left shift & byte offsets into \code{Bytes}; an index that is the bare counter \\
    invariant part & a variable bound before the loop, or a loop-invariant stride & an invariant product inlined into the address \\
    GC-freedom & no allocation or non-inlined call; polls moved to exits by counted loops & loops that call functions or box values \\
    where time goes & inner loops dominated by address arithmetic & arithmetic reached through calls; compute-bound kernels \\
    registers & free registers for the carried pointers & loops at the register limit \\
    \bottomrule
  \end{tabularx}
  \caption{When the Cmm loop passes apply.}
  \label{tab:applicability}
\end{table}

A vectorised ChaCha20, or number-theoretic transforms indexing
\code{a.(j + len)}, have the shapes the passes handle.
